Danh mục này tổng hợp các từ viết tắt, ký hiệu mã và một số thuật ngữ kỹ thuật được sử dụng trong báo cáo. Các mục được lựa chọn theo nội dung thực tế của đề tài website tập luyện thể dục, nhằm giúp người đọc thuận tiện hơn khi theo dõi phần phân tích, thiết kế, triển khai và kiểm thử hệ thống.

\begin{longtable}{|P{0.19\textwidth}|P{0.34\textwidth}|P{0.39\textwidth}|}
\caption*{\textbf{Bảng 1: Danh mục từ viết tắt và thuật ngữ dùng trong báo cáo}}\\
\hline
\textbf{Viết tắt/thuật ngữ} & \textbf{Nội dung đầy đủ} & \textbf{Ghi chú} \\
\hline
\endfirsthead
\caption*{\textbf{Bảng 1: Danh mục từ viết tắt và thuật ngữ dùng trong báo cáo (tiếp theo)}}\\
\hline
\textbf{Viết tắt/thuật ngữ} & \textbf{Nội dung đầy đủ} & \textbf{Ghi chú} \\
\hline
\endhead
API & Application Programming Interface & Giao diện lập trình ứng dụng, dùng để trao đổi dữ liệu giữa giao diện, route xử lý và dịch vụ ngoài. \\
\hline
Auth0 & Auth0 Authentication Platform & Dịch vụ xác thực tài khoản, được dùng trong luồng đăng nhập và quản lý phiên người dùng. \\
\hline
CSS & Cascading Style Sheets & Ngôn ngữ định kiểu giao diện web; trong đề tài có liên quan đến thiết kế hiển thị và bố cục trang. \\
\hline
CRUD & Create, Read, Update, Delete & Nhóm thao tác thêm, xem, sửa và xóa dữ liệu, thường dùng trong chức năng quản trị người dùng và bài tập. \\
\hline
DB & Database & Cơ sở dữ liệu lưu thông tin tài khoản, hồ sơ sức khỏe, bài tập, kế hoạch và tiến trình luyện tập. \\
\hline
Dify/Coze & Dify/Coze Chatbot Platform & Nền tảng chatbot được nhúng hoặc tích hợp để hỗ trợ hỏi đáp, hướng dẫn sử dụng và tư vấn tham khảo. \\
\hline
HTML & HyperText Markup Language & Ngôn ngữ đánh dấu dùng để xây dựng cấu trúc nội dung trên trang web. \\
\hline
HTTP & HyperText Transfer Protocol & Giao thức truyền tải dữ liệu giữa trình duyệt và máy chủ web. \\
\hline
JSON & JavaScript Object Notation & Định dạng dữ liệu có cấu trúc, thường dùng khi route hoặc API trả dữ liệu cho giao diện. \\
\hline
JSX & JavaScript XML & Cú pháp mở rộng của JavaScript, được dùng trong các file giao diện và route của project React/Remix. \\
\hline
NextUI & NextUI Component Library & Thư viện giao diện hỗ trợ xây dựng các thành phần như nút, bảng, form và bố cục hiển thị. \\
\hline
ORM & Object-Relational Mapping & Cơ chế ánh xạ giữa mã nguồn và cơ sở dữ liệu quan hệ, giúp thao tác dữ liệu thông qua object thay vì viết SQL trực tiếp. \\
\hline
PostgreSQL & PostgreSQL Database & Hệ quản trị cơ sở dữ liệu quan hệ được sử dụng để lưu dữ liệu nghiệp vụ của hệ thống. \\
\hline
Prisma & Prisma ORM & Công cụ ORM dùng để kết nối và truy vấn cơ sở dữ liệu trong mã nguồn. \\
\hline
RapidAPI & RapidAPI Platform & Nền tảng cung cấp API bên ngoài, có thể dùng để lấy hoặc tham khảo dữ liệu bài tập/video liên quan. \\
\hline
Recharts & Recharts Chart Library & Thư viện biểu đồ dùng để trực quan hóa dữ liệu thống kê tiến trình luyện tập. \\
\hline
Remix & Remix Web Framework & Framework full-stack dùng để xây dựng route, loader, action và giao diện của website. \\
\hline
SQL & Structured Query Language & Ngôn ngữ truy vấn cơ sở dữ liệu quan hệ; trong project được thao tác chủ yếu thông qua Prisma. \\
\hline
Tailwind CSS & Tailwind Cascading Style Sheets & Framework CSS tiện ích, hỗ trợ định dạng giao diện nhanh và nhất quán. \\
\hline
UI & User Interface & Giao diện người dùng, gồm các màn hình, nút, bảng, form và thành phần tương tác. \\
\hline
URL & Uniform Resource Locator & Đường dẫn tài nguyên hoặc endpoint được dùng khi truy cập trang, API, ảnh hoặc video. \\
\hline
UX & User Experience & Trải nghiệm người dùng khi thao tác với hệ thống, bao gồm mức độ dễ dùng, rõ ràng và nhất quán. \\
\hline
\end{longtable}

\begin{longtable}{|P{0.19\textwidth}|P{0.34\textwidth}|P{0.39\textwidth}|}
\caption*{\textbf{Bảng 2: Quy ước mã yêu cầu, Use Case và kiểm thử}}\\
\hline
\textbf{Ký hiệu mã} & \textbf{Nội dung đầy đủ} & \textbf{Ghi chú} \\
\hline
\endfirsthead
\caption*{\textbf{Bảng 2: Quy ước mã yêu cầu, Use Case và kiểm thử (tiếp theo)}}\\
\hline
\textbf{Ký hiệu mã} & \textbf{Nội dung đầy đủ} & \textbf{Ghi chú} \\
\hline
\endhead
FR & Functional Requirement & Yêu cầu chức năng, ví dụ FR-01, FR-02, dùng để mô tả chức năng hệ thống cần đáp ứng. \\
\hline
NFR & Non-Functional Requirement & Yêu cầu phi chức năng, ví dụ NFR-01, NFR-02, dùng để mô tả hiệu năng, bảo mật, dễ sử dụng và khả năng mở rộng. \\
\hline
UC & Use Case & Ca sử dụng, ví dụ UC-01, UC-02, dùng để mô tả chức năng theo góc nhìn tác nhân. \\
\hline
TC & Test Case & Kịch bản kiểm thử, ví dụ TC-01, TC-02, dùng để kiểm tra luồng chức năng chính của hệ thống. \\
\hline
AUTH & Authentication & Nhóm kiểm thử hoặc chức năng liên quan đến đăng nhập, xác thực và quản lý phiên. \\
\hline
HEALTH & Health Status & Nhóm chức năng hoặc kịch bản liên quan đến hồ sơ sức khỏe người dùng. \\
\hline
EXERCISE & Exercise Management & Nhóm chức năng hoặc kịch bản liên quan đến quản lý và hiển thị bài tập. \\
\hline
REC & Recommendation & Nhóm chức năng hoặc kịch bản liên quan đến gợi ý bài tập theo hồ sơ cá nhân. \\
\hline
PLAN & Workout Plan & Nhóm chức năng hoặc kịch bản liên quan đến lập và quản lý kế hoạch tập luyện. \\
\hline
PROG & User Progress & Nhóm chức năng hoặc kịch bản liên quan đến ghi nhận tiến trình tập luyện. \\
\hline
STAT & Statistics & Nhóm chức năng hoặc kịch bản liên quan đến thống kê và trực quan hóa kết quả luyện tập. \\
\hline
ADMIN & Administrator & Nhóm chức năng hoặc kịch bản liên quan đến quản trị người dùng và dữ liệu bài tập. \\
\hline
\end{longtable}
